\chapter{Постановка задачи}\label{ch:4}

Структурный системный анализ объекта исследования показал, что значительная доля заявок клиентов поступает в виде фотографий, а процесс обработки обращений имеет множество недостатков, таких как длительные ожидания и согласования, просчёты новых сотрудников. Устойчивость и воспроизводимость результатов предварительной оценки подтверждает возможность и необходимость внедрения автоматизированного сервиса оценки повреждений автомобилей по фотографии.

Цель работы - повышение эффективности процесса оценки стоимости ремонта поврежденных автомобилей. Был сформулирован список задач, которые необходимо выполнить для достижения цели:

\begin{enumerate}[label=\arabic*.]
    \item Сформировать датасет для обучения модели путём сбора и фильтрации изображений повреждённых автомобилей и их дальнейшей разметки. 
    \item Разработать модули для обработки и аугментации изображений, устранения недостатков съёмки
    \item Обучить ML-модель обнаружения и классификации деталей и повреждений на основе размеченных данных
    \item Реализовать бизнес-логику оценки стоимости и длительности ремонта
    \item Спроектировать и реализовать микросервисную архитектуру приложения, описать и произвести процесс интеграции с ботом в мессенджере
    \item Описать и разработать архитектуру базы данных
    \item Интеграция CI/CD и процесса поставки программного обеспечения
    \item Провести тестирование работоспособности системы
\end{enumerate}

Предполагаемым результатом работы является реализованный сервис с внедрённой в него ML-моделью распознавания и классификации повреждений автомобиля. 

Эффективность решения можно будет оценить двум основным факторам:
\begin{enumerate}[label=\arabic*.]
    \item Производительность системы
    \begin{itemize}
        \item Время обработки одного изображения
        \item Пропускная способность сервисов (RPS) 
    \end{itemize}
    \item Минимизация человеческого участия - на сколько уменьшится время обработки обращений клиентов
\end{enumerate}